\chapter{Writing Tools}
\label{chap:writing}

\section{Citations}

Entries live in \texttt{references.bib} and are cited by key. A book
\cite{knuth1984texbook}, a journal article \cite{knuth1984literate}, a
conference paper \cite{vaswani2017attention}, a website \cite{latexworkshop}, or
several at once \cite{lamport1994latex,knuth1984texbook}. Only cited entries
appear in the bibliography, numbered in order of appearance.

\section{Acronyms}

Acronyms are declared once in \texttt{acronyms.tex}. The first time,
\verb|\ac{API}| prints the long form, \ac{API}; after that it prints the
short form, \ac{API}. Plurals work with \verb|\acp|: several \acp{IDE}.
Every acronym used appears in the list at the beginning of the report.

\section{Footnotes and quotes}

Footnotes are numbered automatically.\footnote{And placed at the bottom of
the page.} Short quotes use \verb|\enquote|: \enquote{the quotation marks
follow the document language}. Long quotes get their own block:

\begin{displayquote}
  A long quotation is set apart from the surrounding text and indented on
  both sides, so the reader sees at a glance where it starts and ends.
\end{displayquote}

\section{Colored boxes}

\begin{notebox}
  Boxes made with \texttt{tcolorbox} highlight information the reader should
  not miss. Their style is defined once in \texttt{main.tex}.
\end{notebox}

\begin{warningbox}[Careful]
  The title of a box can be changed for each use.
\end{warningbox}

\section{Lists}

\begin{itemize}
  \item A bullet point, with \textbf{bold}, \emph{emphasis} and \underline{underline}.
  \item A nested list:
  \begin{itemize}
    \item second level,
    \item still second level.
  \end{itemize}
\end{itemize}

\begin{enumerate}[label=\textbf{Step \arabic*.}, leftmargin=*]
  \item Custom labels come from the \texttt{enumitem} package.
  \item The numbering continues by itself.
\end{enumerate}

\begin{description}
  \item[latexmk] runs the right tools as many times as needed.
  \item[SyncTeX] links each line of the source to its place in the \ac{PDF}.
\end{description}

\section{Cross references and links}

Labels make every number clickable: \cref{fig:tikz}, \cref{tab:merged},
\cref{lst:python}, \cref{eq:mean}, \cref{def:precision} and
\cref{chap:figures}. Page numbers work too: \cref{tab:long} starts on
page~\pageref{tab:long}. Web links can be written as is,
\url{https://github.com/Ghorbel37/LocalLeaf}, or behind text, like
\href{https://www.overleaf.com/learn}{Overleaf's guides}.
